\documentclass[10pt,a4paper]{amsart}
\usepackage[margin=19mm]{geometry}
\usepackage{amsmath,amssymb,mathtools}
\usepackage[scr=rsfs,cal=euler]{mathalfa}
\usepackage{xfrac}
\usepackage[unicode,hidelinks,bookmarks=false]{hyperref}
\newcommand{\ie}{i.e.\ }
\newcommand{\cf}{cf.\ }
\newcommand{\resp}{respectively\ }
\newcommand{\Subsection}{\subsection}
\providecommand{\nobreakdash}{\nobreak\hbox{-}}
\setlength{\emergencystretch}{2em}
\multlinegap0pt
\usepackage{float}
\usepackage{amssymb}
\usepackage{braket}
\usepackage{mathtools}
\usepackage{tikz}
\usepackage{xcolor}
\newtheorem{proposition}{Proposition}
\newtheorem{lemma}{Lemma}
\newcommand{\C}{\mathbb{C}}
\newcommand{\N}{\mathbb{N}}
\newcommand{\R}{\mathbb{R}}
\newcommand{\Z}{\mathbb{Z}}
\newcommand{\e}{\operatorname{e}}
\title{Effective estimates for exponential sums with multiplicative coefficients}
\author{Nicolas Robles}
\date{Source: arXiv:2609.11070v1 (CC BY 4.0)}
\begin{document}
\maketitle

\noindent\emph{Source and licence.} Effective estimates for exponential sums with multiplicative coefficients. Version arXiv:2609.11070v1 (CC BY 4.0), distributed by arXiv under the Creative Commons Attribution 4.0 International licence, http://creativecommons.org/licenses/by/4.0/, as recorded in the arXiv licence metadata for this exact version and shown on the abstract page https://arxiv.org/abs/2609.11070v1. Full licence notice: \texttt{licenses/arxiv-2609.11070-CC-BY-4.0.txt}.\par\smallskip

\noindent\textit{Editorial scope.} Counted unit: the article's proposition prop:effective-maximal-prime (source0.tex lines 1161--1179). The article's complete original proof of it is delivered with it (source0.tex lines 1181--1265, one \texttt{proof} environment of 85 lines). No mathematics is written, repaired or re-derived: the statement and the proof are the article's own lines, in the article's own notation, and no retained line is substituted or re-flowed.  Retained uncounted as the source-local context the proof needs: lines 608--616; lines 898--900; lines 902--905; lines 985--993; lines 1010--1018.  Nothing was omitted from any retained span, so the package carries no omission record.  No retained line contains a citation, so the wrapper carries no bibliography and no bibliography entry.  No retained line contains a figure environment, an image-inclusion command or drawing code, so no image is delivered and the package declares no asset dependency.  Editorial formatting only, in addition: an AMS \texttt{amsart} preamble that restates the article's own declarations and macros used by the retained lines, namely the environments \texttt{proposition}, \texttt{lemma} on separate counters, together with the standard packages those lines need.  The retained environments print in this document as Proposition 1, Lemma 1 in order of appearance, whereas the article prints the same items with the numbering of its own full document.  The complete native source of the article as distributed in the arXiv e-print for this exact version is bundled unchanged as \texttt{sources/209984/source0.tex}.
\begin{proposition}[Brun-Titchmarsh on an arbitrary interval]\label{prop:BT}
There is an effective absolute constant $C_\mathrm{BT}$ such that, for
real $x\ge h\ge r\ge1$, integers $r,b$ with $(b,r)=1$,
\begin{equation}\label{eq:BT}
 \#\{p:x-h<p\le x,\ p\equiv b\pmod r\}
 \le C_\mathrm{BT}\frac{h}{\varphi(r)\log(2h/r)}.
\end{equation}
The constant is independent of the location $x$ of the interval.
\end{proposition}

\begin{equation}\label{eq:maximal-all-residues}
 \sum_{b\bmod r}M_c(\theta+b/r)^2\ll (U+r)E(c).
\end{equation}

\begin{equation}\label{eq:maximal-reduced-residues}
 \frac1{\varphi(r)}\sum_{b\in\mathcal R(r)}M_c(\theta+b/r)^2
 \ll\log\log(3U)E(c).
\end{equation}

\begin{lemma}\label{lem:covering}
For $P>0$, $c\in\C^U$, $a\in\Z$, $r\in\N$, $(a,r)=1$, and
$\beta\ne0$,
\begin{equation}\label{eq:covering}
 \sum_{b\bmod r}\int_P^{2P}M_c(ab/r+\beta t)^2\,dt
 \ll (rP+|\beta|^{-1})E(c).
\end{equation}
The constant is absolute and effective.
\end{lemma}

\begin{lemma}\label{lem:variation}
Let $I$ be a compact interval of length $h>0$ and let $F:I\to\C$
be continuously differentiable. Then
\begin{equation}\label{eq:interval-variation}
 \sup_{t\in I}|F(t)|^2
 \le\frac2h\int_I|F(t)|^2\,dt
       +2h\int_I|F'(t)|^2\,dt.
\end{equation}
\end{lemma}

\begin{proposition}\label{prop:effective-maximal-prime}
Let $N>P$, $U=\lfloor N/P\rfloor$, and
$\alpha=a/r+\beta$, where $a\in\Z$, $r\in\N$, $(a,r)=1$, and
$\beta\in\R$. Set $B=\max\{1,N|\beta|\}$ and $D=rB$.
For $c\in\C^U$, write $E=E(c)$ and
\begin{equation}\label{eq:Qstar-def}
 Q_P^*(c)=\sum_{P<p\le2P}M_c(\alpha p)^2\log p.
\end{equation}
If $P\ge16D^2$, then
\begin{equation}\label{eq:effective-prime-general}
 Q_P^*(c)\ll\frac{rP+N/B}{\varphi(r)}E.
\end{equation}
If also $U\le r$, then
\begin{equation}\label{eq:effective-prime-short}
 Q_P^*(c)\ll P\log\log(3U)E.
\end{equation}
The constants are absolute and effective. There is no multiplicativity
hypothesis on $c$.
\end{proposition}

\begin{proof}
Put $J=\lceil B\rceil$ and $h=P/J$. Partition $(P,2P]$ into the
half-open intervals
$I_j=(P+(j-1)h,P+jh]$, $1\le j\le J$.
Their closures are used when taking suprema and integrals. Since
$B\le J\le2B$, $PU\le N$, and $N|\beta|\le B$,
\begin{equation}\label{eq:interval-parameters}
 h\ge\frac{P}{2B},\qquad |\beta|hU\le1,\qquad
 \frac h r\ge\frac{P}{2D}\ge2\sqrt P.
\end{equation}
Also $P\ge16D^2\ge16r^2>r$, so every prime in $(P,2P]$ is
coprime to $r$.

For $b\in\mathcal R(r)$, Proposition~\ref{prop:BT} applies to $I_j$
with upper endpoint $P+jh$ and length $h$. It gives
\begin{equation}\label{eq:effective-prime-mass}
 \sum_{\substack{p\in I_j\\p\equiv b\pmod r}}\log p
 \le C_\mathrm{BT}\frac{h\log(2P)}{\varphi(r)\log(2h/r)}
 \ll\frac h{\varphi(r)}.
\end{equation}
The last constant is effective: by \eqref{eq:interval-parameters},
$\log(2h/r)\ge\log(4\sqrt P)$, whose ratio to $\log(2P)$ is bounded
below by a positive absolute constant.

For $0\le K\le U$ put
\[
 F_{K,b}(t)=\sum_{n\le K}c_n\e\bigl(n(ab/r+\beta t)\bigr),
 \qquad d_n=nc_n\quad(1\le n\le U).
\]
Then $E(d)\le U^2E$ and
$|F'_{K,b}(t)|\le2\pi|\beta|M_d(ab/r+\beta t)$.
Apply Lemma~\ref{lem:variation} to each fixed $K$ on the closure of
$I_j$. Bounding its two integrands by the corresponding maximal
expressions gives a bound independent of $K$, and hence
\[
 \sup_{t\in I_j}M_c(ab/r+\beta t)^2
 \le\frac2h\int_{I_j}M_c(ab/r+\beta t)^2\,dt
 +8\pi^2h|\beta|^2\int_{I_j}M_d(ab/r+\beta t)^2\,dt.
\]
For a prime $p\equiv b\pmod r$, periodicity gives
$M_c(\alpha p)=M_c(ab/r+\beta p)$.
Multiply the last inequality by the prime-mass bound
\eqref{eq:effective-prime-mass}, and sum over $j$ and $b$.
We obtain the positive transfer estimate
\begin{equation}\label{eq:positive-sieve-transfer}
 Q_P^*(c)\ll\frac1{\varphi(r)}\sum_{b\in\mathcal R(r)}
 \int_P^{2P}\big(M_c(ab/r+\beta t)^2
   +h^2|\beta|^2 M_d(ab/r+\beta t)^2\big)\,dt.
\end{equation}
Only the fixed prefix polynomials were differentiated, not the maximum
and not any prime-dependent truncation.

If $B>1$, enlarge the residue sum to all classes and apply
Lemma~\ref{lem:covering}, once to $c$ and once to $d$.
Equation~\eqref{eq:positive-sieve-transfer} is then
\[
 \ll\frac{rP+|\beta|^{-1}}{\varphi(r)}
                  (E+h^2|\beta|^2E(d))
 \le\frac{2(rP+N/B)}{\varphi(r)}E
\]
up to an effective absolute constant. Here
$|\beta|^{-1}=N/B$ and \eqref{eq:interval-parameters} was used.
This proves \eqref{eq:effective-prime-general} in this case.

If $B=1$, including $\beta=0$, apply
\eqref{eq:maximal-all-residues} at each fixed $t$, again to both vectors.
It bounds \eqref{eq:positive-sieve-transfer} by
\[
 \ll\frac{P(U+r)}{\varphi(r)}
                  (E+h^2|\beta|^2E(d))
 \ll\frac{N+rP}{\varphi(r)}E.
\]
This is \eqref{eq:effective-prime-general} with $B=1$.

Finally, when $U\le r$, keep reduced residues in
\eqref{eq:positive-sieve-transfer} and apply
\eqref{eq:maximal-reduced-residues} at each $t$.
The outcome is
\[
 Q_P^*(c)\ll P\log\log(3U)
                  (E+h^2|\beta|^2E(d))
 \ll P\log\log(3U)E.
\]
This proves \eqref{eq:effective-prime-short} for every $\beta$, including zero.
\end{proof}

\end{document}
